\chapter{多变量函数的微分学}
\section{多变量函数及其连续性}


\subsection{点集的结构}
多变量函数是定义在高维空间点集上的函数，为了满足所要研究函数的性质，需要对这个点集所具有的结构做进一步的规定.

\begin{definition}
    定义集合$\mathbb{R}^n$是有序$n$元实数组的集合：
    \[\mathbb{R}^n \coloneq \{(x_1,x_2,\cdots,x_n)\,|\,x_i \in \mathbb{R}\}\]
    其中的元素$\vb{x} = (x_1,x_2,\cdots,x_n)$称为$\mathbb{R}^n$中的一个点.
\end{definition}

\subsubsection{实线性空间}
\begin{definition}
    对于$\mathbb{R}^n$中的两点$\vb{x}=(x_1,\cdots,x_n),\vb{y}=(y_1,\cdots,y_n)$，定义加法运算：
    \[\vb{x} + \vb{y} \coloneq (x_1+y_1,\cdots,x_n+y_n)\]
    定义数乘运算：
    \[\forall \lambda \in \mathbb{R},\quad \lambda\vb{x} \coloneq (\lambda x_1,\cdots,\lambda x_n)\]
    定义有上述运算的空间称为\overtext{实线性空间}{real linear space}或\overtext{实向量空间}{real vector space}.
\end{definition}


\subsubsection{赋范空间}
\begin{definition}
    在线性空间上定义\overtext{范数}{norm}运算：
    \[\left\|\vb{x}\right\| \coloneq \sqrt{x_1^2 + \cdots + x_n^2}\]
    称定义有范数的线性空间为\overtext{赋范空间}{normed space}\!\wiki{https://zh.wikipedia.org/wiki/\%E8\%B3\%A6\%E7\%AF\%84\%E5\%90\%91\%E9\%87\%8F\%E7\%A9\%BA\%E9\%96\%93}.
\end{definition}


\subsubsection{度量空间}
\begin{definition}
    在赋范空间上定义\overtext{距离}{metric}的概念：
    \[\rho(\vb{x},\vb{y}) \coloneq \left\|\vb{x} - \vb{y}\right\| = \sqrt{\sum_{i=1}^{n}(x_i-y_i)^2}\]
    其满足距离的四条公理：
    \begin{enumerate}[circ]
        \item \textbf{非负性：}$\rho(\vb{x,y}) \geqslant 0$.
        \item \textbf{不可区分性}$\rho(\vb{x,y})\Leftrightarrow \vb{x=y}$.
        \item \textbf{对称性：}$\rho(\vb{x,y}) = \rho(\vb{y,x})$.
        \item \textbf{三角不等式：}$\rho(\vb{x,y})\leqslant\rho(\vb{x,z}) + \rho(\vb{y,z})$.
    \end{enumerate}
    称定义有距离的空间为\overtext{度量空间}{metric space}\!\wiki{https://zh.wikipedia.org/wiki/\%E5\%BA\%A6\%E9\%87\%8F\%E7\%A9\%BA\%E9\%97\%B4}.
\end{definition}


\subsubsection{拓扑空间}
为了定义极限，要为集合赋予「一点附近」的概念.这样的空间称为\overtext{拓扑空间}{topological space}\!\wiki{https://zh.wikipedia.org/wiki/\%E6\%8B\%93\%E6\%89\%91\%E7\%A9\%BA\%E9\%97\%B4}.在拓扑空间中，可以定义以下概念：

\begin{definition}
    设有一点$\vb{x}_0$，称满足如下定义的集合为一个\overtext{开球}{open ball}：
    \[B(\vb{x}_0,r) = \{\vb{x}\,|\,\rho(\vb{x},\vb{x}_0) < r\}\]
    除去中心点$\vb{x}_0$，则得到\overtext{去心开球}{punctured open ball}的定义：
    \[B_-(\vb{x}_0,r) = \{\vb{x}\,|\,0 < \rho(\vb{x},\vb{x}_0) < r\}\]
\end{definition}

\begin{definition}
    邻域是包含某个开球的集合：
    \begin{itemize}
        \item 称$U$是$\vb{x}_0$的\overtext{邻域}{neighborhood}，当且仅当$\exists r>0$使得$B(\vb{x}_0,r)\subseteq U$.
        \item 称$U_-$是$\vb{x}_0$的\overtext{去心邻域}{punctured neighborhood}，当且仅当$\exists r>0$使得$B_-(\vb{x}_0,r)\subseteq U_-$.
    \end{itemize}
\end{definition}

从邻域的定义不难发现，邻域只要能够包含某个开球即可，并没有固定的形状，常用的是球形或方形.

\begin{definition}
    若对于一个点集$E$，$\exists R>0$使得$E \subset B(O,R)$，则称$E$是\textbf{有界的}，其中$O$是坐标原点.

    定义有界非空点集$E$的\overtext{直径}{diameter}为集合中任意两点距离的上确界，记作：
    \[\mathrm{diam}\, E = \sup\{\rho({\vb{x,y}})\,|\,\vb{x,y}\in E\}\]
\end{definition}

\begin{definition}
    设$E \subset \mathbb{R}^n$是一个点集，$\mathbb{R}^n$中不属于$E$的部分称为$E$的\overtext{余集}{complement}，记作$E^c$，显然有$(E^c)^c = E$.平面上的点按照其与$E$的关系分为以下三种：
    \begin{enumerate}[circ]
        \item 如果$\exists r>0$，使得$B(\vb{x}_0,r) \subset E$，则称$\vb{x}_0$是$E$的\overtext{内点}{interior point}.$E$中全体内点组成的集合称为$E$的\overtext{核}{interior}，记作$E^\circ$.
        \item 如果$\exists r>0$，使得$B(\vb{x}_0,r) \subset E^c$，则称$\vb{x}_0$是$E$的\overtext{外点}{exterior point}.
        \item 如果$\forall r>0$，使得$B(\vb{x}_0,r)$中既有$E$的点，又有$E^c$的点，则称$\vb{x}_0$为$E$的\overtext{边界点}{boundary point}.$E$的全体边界点的集合称为$E$的\overtext{边界}{boundary}，记作$\partial E$.
    \end{enumerate}
\end{definition}

显然按照边界点的定义，边界点既可以在$E$中，也可以不在$E$中.

\begin{definition}
    \begin{itemize}
        \item 对于$E$中的点，如果$\exists r>0$，使得$B(\vb{x}_0,r) \cap E = \vb{x}_0$，那么称$\vb{x}_0$是$E$的\overtext{孤立点}{}.
        \item 对于空间中一点$\vb{x}_0$，如果$\forall r > 0$，使得$B_-(\vb{x}_0,r)$中总是有$E$中的点，称$\vb{x}_0$为$E$的\overtext{聚点}{}.
    \end{itemize}
\end{definition}

根据上述定义，$E$的聚点不一定在$E$中，$E$的边界点要么是聚点，要么是孤立点.不孤立的$E$的边界点则一定是聚点，在$E$中的内点也一定是聚点.

\begin{definition}
    设$E \subset \mathbb{R}^n$，若$E = E^\circ$，即$E$中的点都是$E$的内点，则称$E$是一个\overtext{开集}{open set}；若$E$的余集$E^c$是开集，则称$E$是\overtext{闭集}{close set}.
\end{definition}

按照上述定义，在$\mathbb{R}^n$中，只有整个空间$\mathbb{R}^n$和空集$\emptyset$既是\overtext{既开又闭集}{clopen set}.一方面，它们中没有不是内点的点，另一方面，二者互为余集，因此也是闭集.

\begin{definition}
    设$E \in \mathbb{R}^n$，记$\overline{E} = E \cup \{E\text{的所有聚点}\}$，称为$E$的\overtext{闭包}{}.
\end{definition}
可以认为$E$的闭包就是$E$本身加上所有从极限上能够靠近到的点，或者说是$E$中所有点加上不在$E$中的$E$的边界点.

开集和闭集具有如下性质：
\begin{enumerate}[property]
    \item 两个开集的并或交仍是开集，两个闭集的并或交仍是闭集.
    \begin{proof}
        以开集之交的证明为例：设有开集$E_1,E_2$，$\forall M \in E_1 \cap E_2$，满足$M \in E_1$且$M \in E_2$.对$M \in E_1$，$\exists r_1 > 0$使得$B(M,r_1) \subset E_1$；对$M \in E_2$，$\exists r_2 > 0$使得$B(M,r_2) \subset E_2$.

        不妨取$r = \min\{r_1,r_2\}$，则对$\forall M \in E_1 \cap E_2$，$\exists r$满足$B(M,r) \subset (E_1 \cap E_2)$.
    \end{proof}
    \item $E\subset\mathbb{R}^n$是开集，当且仅当$\partial E \cap E = \emptyset$.
    \item $E\subset\mathbb{R}^n$是闭集，当且仅当$\partial E \subset E$.
    \item $E\subset\mathbb{R}^n$是闭集，当且仅当$E = \overline{E}$.
\end{enumerate}
\begin{mistake}
    设有命题$p,q,q'$，若有$p\Leftrightarrow q$且$p\Leftrightarrow q'$，则命题$q$和$q'$在表述上可以不一致，但是在逻辑上一定等价，即$q \Leftrightarrow q'$.
\end{mistake}

\begin{definition}
    对于点集$E \subset \mathbb{R}^n$，如果存在两个非空开集$A,B$使得$E = A \cup B$时，则称$E$是\textbf{不}\!\!\overtext{连通}{connected}的，否则称$E$是连通的.
\end{definition}


\begin{definition}
    \textbf{连续曲线}是由如下参数方程表示函数的图像：
    \[\vb{r}(t) = (x(t),t(t)),\quad t\in[\alpha,\beta]\]
    其中$x(t),y(t)$均为$[\alpha,\beta]$上的连续函数.若上式满足对任意的$\alpha \leqslant t_1 < t_2 \leqslant \beta$有$(x(t_1),y(t_1)) \neq (x(t_2),y(t_2))$，即该曲线不存在自相交，则称曲线为一条\textbf{简单曲线(Jordan曲线)}.

    若曲线的两端点重合，即$(x(\alpha),y(\alpha)) = (x(\beta),y(\beta))$，则称为\textbf{闭曲线}.
\end{definition}

\begin{definition}
    对于点集$E$中任意两点$P,Q$，若在$E$中存在一条连续曲线将$P$和$Q$连接起来，则称$E$是\overtext{道路连通}{path connected}或\overtext{弧连通}{arc connected}的.
\end{definition}

\begin{theorem}
    道理连通集一定是连通的，反之则不然.
\end{theorem}

\begin{proof}
    取$E = A \cup B$且$A \cap B = \emptyset$.取$\forall M_1 \in A,\forall M_2 \in B$，因为$E$是道路连通的，因此存在连续曲线：
    \[\vb{r}(t) = (x(t),y(t)),\quad t\in[\alpha,\beta]\quad\text{且}\quad\vb{r}(\alpha) = M_1,\vb{r}(\beta) = M_2\]

    设$t_0 = \inf\{t\,|\,\vb{r}(t) \in B\}$，此时分两种情况讨论.若$\vb{r}(t_0)\in A$，则$\forall\varepsilon > 0$，$\exists\delta>0$使得在$|t-t_0|<\delta$时有：
    \[|x(t) - x(t_0)| < \frac{\varepsilon}{2},\qquad |y(t) - y(t_0)| < \frac{\varepsilon}{2}\]
    若是取$t=t',\quad t_0 < t < t_0 + \delta$，则显然$\vb{r}(t') \in B$.

    即在$B(\vb{r}(t_0),\varepsilon)$中对$\forall \varepsilon>0$都存在$B$中的点，则表明$\vb{r}(t_0)\in A$是$B$的聚点，满足连通的定义.
\end{proof}

\begin{definition}
    道路连通的开集称为\overtext{开区域}{open domain}，开区域的闭包称为\overtext{闭区域}{close region}，统称为\overtext{区域}{region}.
\end{definition}

\begin{definition}
    如果区域$D$内任意一条简单闭曲线的内部仍在$D$内，即$D$内没有空洞，那么称$D$是\overtext{单连通}{simply connected}的，反之则称为\overtext{多连通}{multiply connected}的.
\end{definition}


\subsection{平面上的点集}
\subsubsection{平面点列的极限}
\begin{definition}
    设$\{M_n\}$是平面点列，如有$M_0 \in \mathbb{R}^2$使得$\lim \rho(M_n,M_0) = 0$，则称点列$\{M_n\}$为\textbf{收敛点列}，$M_0$称为点列的\textbf{极限}.记作：
    \[\lim_{n \to \infty} M_n  = M_0\]
\end{definition}

平面上的点列不一定需要以直线逼近其极限，可以由任意的轨迹逼近.

\begin{theorem}
    设$M_n(x_n,y_n)$，则有：
    \[\lim M_n = M_0 \quad\Leftrightarrow\quad \begin{cases}
        \lim x_n = x_0 \\
        \lim y_n = y_0
    \end{cases}\]
\end{theorem}

\begin{proof}
    注意到有以下不等式：
    \begin{gather*}
        |x_n-x_0| \leqslant \rho(M_n,M_0) \leqslant |x_n-x_0| + |y_n-y_0| \\
        |y_n-y_0| \leqslant \rho(M_n,M_0) \leqslant |x_n-x_0| + |y_n-y_0|
    \end{gather*}
    根据三明治法则，即得证.
\end{proof}

\begin{theorem}
    平面上的有界点列必有收敛子列.
\end{theorem}

\begin{proof}
    设有$M_n(x_n,y_n)$，若$\{x_n\}$有界，则必有收敛子列$\{x_{n_k}\}$，则$\{x_{n_k}\}$的子列$\{x_{n_k'}\}$也收敛.

    若$\{y_n\}$有界，则必有子列$\{y_{n_k}\}$有界，则其进一步有有界子列$\{y_{n_k'}\}$收敛.

    则平面上点列$\{M_n\}$必有收敛子列$\{M_{n_k'}\}$.
\end{proof}


\subsection{多变量函数}
\subsubsection{多变量函数的定义}
$n$维实数组空间$\mathbb{R}^n = \{\vb{x} = (x_1,x_2,\cdots,x_n)\,|\,x_i\in\mathbb{R},i=1,2,\cdots,n\}$中的点$\vb{x}$对应一个向量：
\[\vb{x} = x_1\vb{e}_1 + x_2\vb{e}_2 + \cdots + x_n\vb{e}_n\]
其中$\vb{e}_i$是$\mathbb{R}_n$的标准正交基.

\begin{definition}
    设$D \subset \mathbb{R}^n$，则从$D$到直线上的映射称为一个\overtext{$n$元函数}{$n$-variable function}，记作：
    \[f:D \to \mathbb{R} \quad\text{或}\quad z = f(x_1,x_2,\cdots,x_n)\]
\end{definition}

对于二元函数$z=f(x,y),(x,y)\in D \subset \mathbb{R}^n$，如果把$z$所在的数轴垂直于$Oxy$平面放置并形成右手系，则得到一个三维空间$Oxyz$上的一张曲面：
\[S = \{(x,y,f(x,y))\,|\,(x,y)\in D\}\]
若在$Oxy$平面上取使得$f(x,y)$等于一固定常数$c$的点的集合：
\[L_c = \{(x,y)\,|\,f(x,y)=c\}\]
则得到了$Oxy$平面上的一条或几条曲线，或若干点，或为空集，称为\overtext{等值线}{contour line}.


\subsection{多变量函数的复合}
将两个二元函数复合，得到的仍是二元函数；将二元函数和一元函数复合得到的则是一元函数.
\begin{figure}[H]
    \centering
    \includegraphics[width=\textwidth]{composition-of-multivar-func/fig.pdf}
    \caption{多变量函数的复合映射过程}
\end{figure}


\subsection{向量值函数}
向量值函数就是一个从$n$维空间$\mathbb{R}^n$到$m$维空间$\mathbb{R}^m$的映射：
\[\vb{f}: D \to \mathbb{R}^m,\quad D \subset \mathbb{R}^n \quad\text{或}\quad \vb{f}:(x_1,x_2,\cdots,x_n)\mapsto(y_1,y_2,\cdots,y_m)\]
应当注意的是，向量值函数$\vb{f}$的每一个分量$y_i$都是所有自变量的多元函数：
\[y_i = f_i(x_1,x_2,\cdots,x_n),\quad i = 1,2,\cdots,m\]
当维数$m=n$时，如果映射可逆，那么称为相同维数空间之间的变换，可逆意味着变换过程中没有发生信息丢失.

常用的向量值函数有$\mathbb{R}\to\mathbb{R}^3$的映射：
\[\vb{r}(t) = x(t)\vb{i} + y(t)\vb{j} + z(t)\vb{k},\quad t\in[\alpha,\beta]\subset\mathbb{R}\]
这将一维的参数$t$映射到了三维空间的一条曲线上，即空间曲线的参数方程表示，称为\overtext{参数曲线}{parametric curve}.

还有$\mathbb{R}^2 \to \mathbb{R}^3$的映射：
\[\vb{r}(u,v) = x(u,v)\vb{i} + y(u,v)\vb{j} + z(u,v)\vb{k},\quad (u,v)\in D \subset \mathbb{R}^2\]
这是将二维参数数组$(u,v)$映射到了三维空间中的一张曲面上，称为\overtext{参数曲面}{parameter surface}.


\subsection{多变量函数的极限}
\begin{definition}
    设$z=f(x,y),(x,y)\in D$，$M_0(x_0,y_0)$是$D$的聚点.若是对$\forall \varepsilon > 0$，$\exists \delta > 0$，当$M(x,y)\in D$满足：
    \begin{gather*}
        \text{圆心去心邻域：}0 < \rho(M,M_0) = \sqrt{(x-x_0)^2 + (y-y_0)^2} < \delta\\
        \text{方形去心邻域：}|x-x_0|<\delta,|y-y_0|<\delta,\quad(x,y)\neq (x_0,y_0)
    \end{gather*}
    使得$|f(M)-a|<\varepsilon$，则称当$M$趋于$M_0$时，$f(M)$以$a$为\textbf{极限}，记作：
    \[\lim_{M \to M_0}f(M) = a \quad\text{或}\quad \lim_{(x,y)\to(x_0,y_0)}f(x,y) = a \quad\text{或}\quad \lim_{\substack{x \to x_0 \\ y \to y_0}}f(x,y) = a\]
\end{definition}

\begin{question}
    证明：$\displaystyle \lim_{(x,y)\to(x_0,y_0)}\frac{x^2y^2}{x^4+y^2} = 0$.
\end{question}

\begin{proof}
    \[\left|\frac{x^2y^2}{x^4+y^2}\right| = \frac{|y|}{|x^4+y^2|}\cdot|x^2 y| \leqslant \frac{|y|}{|x^4+y^2|}\cdot\frac{|x^4+y^2|}{2} = \frac{|y|}{2}\]
    若有极限成立，则上式应满足$\forall \varepsilon > 0$，$\exists\delta>0$使得$|x-0|<\delta,|y-0|<\delta$时有$\frac{|y|}{2} < \varepsilon$.故而取$\delta = 2\varepsilon$，使得上述极限存在且等于$0$.
\end{proof}

\begin{question}
    证明：$\displaystyle \lim_{(x,y)\to(x_0,y_0)}\frac{x^3+y^3}{x^2+y^2} = 0$.
\end{question}

\begin{proof}
    取$x = r\cos\theta, y = r\sin\theta, r \to 0$，有：
    \[\left|\frac{r^3(\cos^3\theta + \sin^3\theta)}{r^2}\right| \leqslant r\]
    对$\forall \varepsilon > 0$，取$\delta = \varepsilon$，使得$r = \rho((x,y),0) < \delta$时有：
    \[\left|\frac{x^3+y^3}{x^2+y^2} - 0\right| \leqslant \varepsilon\]
    故而极限存在且为$0$.
\end{proof}

\begin{theorem}
    设$z=f(x,y),(x,y)\in D$，有$D(x_0,y_0)$是$D$的一个聚点，则有函数极限：
    \[\lim_{(x,y)\to(x_0,y_0)}f(x,y) = a\]
    的充分必要条件是对任何收敛于$(x_0,y_0)$的点列$(x_n,y_n),(x_n,y_n)\neq(x_0,y_0)$都满足：
    \[\lim_{n \to \infty}f(x_n,y_n) = a\]
\end{theorem}

根据\cref{subsec:LimitProperty}中性质7的证明，分别对上述定理中的$x,y$分别证明，而后合并为$\rho((x,y),(x_0,y_0))$进一步证明即可.

\begin{theorem}
    设$z = f(u,v)$，其中$u=u(x,y),v=v(x,y)$，若：
    \begin{enumerate}[circ]
        \item $\lim\limits_{(u,v) \to (u_0,v_0)}f(u,v) = a$
        \vspace{0.5em}
        \item $\lim\limits_{(x,y) \to (x_0,y_0)}u(x,y) = u_0,\quad \lim\limits_{(x,y) \to (x_0,y_0)}v(x,y) = v_0$
    \end{enumerate}
    则$\lim\limits_{(x,y) \to (x_0,y_0)}f(u(x,y),v(x,y)) = a$.
\end{theorem}

\begin{proof}
    对于$\forall \varepsilon > 0$，根据条件\circnum{1}，必定$\exists\delta' > 0$，使得当$|u-u_0|<\delta',|v-v_0|<\delta'$时，有$|f(u,v)-a|<\varepsilon$.

    根据条件\circnum{2}，对于一个给定的$\delta'>0$，一定$\exists \delta>0$在$|x-x_0|<\delta,|y-y_0|<\delta$且$(x,y)\neq(x_0,y_0)$时，满足$|u-u_0|<\delta',|v-v_0|<\delta'$.

    综上所述，对于$\forall \varepsilon>0$，$\exists \delta>0$，在$|x-x_0|<\delta,|y-y_0|<\delta$且$(x,y)\neq(x_0,y_0)$时，满足：
    \[|f(u(x,y),v(x,y))-a|<\varepsilon \quad\Rightarrow\quad \lim_{(x,y) \to (x_0,y_0)}f(u(x,y),v(x,y)) = a\]
\end{proof}

\begin{definition}
    设函数$f(x,y)$在$(x_0,y_0)$附近有定义，先后对$x,y$求极限：
    \[\lim_{y \to y_0} \lim_{x \to x_0}f(x,y),\qquad \lim_{x \to x_0} \lim_{y \to y_0}f(x,y)\]
    如果上述极限存在，那么称之为函数$f(x,y)$在$(x_0,y_0)$处的\overtext{累次极限}{iterated limit}.
\end{definition}

通过举反例不难说明，累次极限与多变量函数的极限没有必然的联系，一个存在不代表另一个也存在.

\begin{theorem}
    设有$\lim\limits_{(x,y) \to (x_0,y_0)}f(x,y) = A$，如果对$y\neq y_0$，有$\varphi(y) = \lim\limits_{x \to x_0}f(x,y)$存在，则：
    \[\lim_{y \to y_0}\left(\lim_{x \to x_0}f(x,y)\right) = A\]
\end{theorem}

\begin{proof}
    对$\forall \varepsilon>0$，$\exists \delta > 0$，在$|x-x_0|<\delta,|y-y_0|<\delta$时满足：
    \[|f(x,y)-A| < \frac{\varepsilon}{2} \quad\Leftrightarrow\quad A - \frac{\varepsilon}{2} < f(x,y) < A + \frac{\varepsilon}{2}\]
    令$x \to x_0$，上式变为：
    \[A - \frac{\varepsilon}{2} \leqslant \varphi(y) \leqslant A - \frac{\varepsilon}{2} \quad\Rightarrow\quad |\varphi(y) - A| \leqslant \frac{\varepsilon}{2} < \varepsilon\]
    即对$\forall \varepsilon>0$，$\exists\delta>0$在$|y-y_0|$时满足$|\varphi(y) - A| < \varepsilon$，即：
    \[\lim_{y \to y_0}\varphi(y) = A \quad\Rightarrow\quad \lim_{y \to y_0}\left(\lim_{x \to x_0}f(x,y)\right) = A\]
    由此得证.
\end{proof}

\begin{theorem}
    设$f(x,y)$在$M_0(x_0,y_0)$附近$B_-(M_0,r)$有定义，且在$M_0$处有有限极限$a$，$x=x(t),y=y(t)$在$0 < |t-t_0| < l$有定义，使得$\{(x(t),y(t))\,|\,0<|t-t_0|<l\}\subset B_-(M_0,r)$且$\lim\limits_{t \to t_0}x(t) = x_0,\lim\limits_{t \to t_0}y(t) = y_0$，则：
    \[\lim_{t \to t_0}f(x(t),y(t)) = a\]
\end{theorem}

\begin{proof}
    已知：
    \[\lim_{(x,y)\to(x_0,y_0)}f(x,y) = a\]
    即$\forall \varepsilon>0$，$\exists \delta_1 > 0$在$|\rho((x,y),(x_0,y_0))|<\delta_1$时满足$|f(x,y)-a|<\varepsilon$.

    现选取一条任意路径$x=x(t),y=y(t)$，在$t\to t_0$时有$x \to x_0,y \to y_0$.即对$\forall \varepsilon > 0$，$\exists \delta_2>0$，使得$0 < \sqrt{(x(t)-x_0)^2+(y(t)-y_0)^2} = \rho((x,y),(x_0,y_0)) < \delta_2$时满足：
    \[|x(t)-x_0|<\varepsilon,\qquad |y(t)-y_0|<\varepsilon\]
    不妨取$\delta = \min\{\delta_1,\delta_2\}$，即这个$\delta$对任意一条路径在$|\rho((x,y),(x_0,y_0))|<\delta$时满足：
    \[\lim_{t \to t_0}f(x(t),y(t)) = a\]
\end{proof}

上述定理表明，只要极限存在，则平面上的点沿任何路径趋于$(x_0,y_0)$时，极限都存在且一致.这个定理主要用于证明极限不存在的情形，只要能找到一条路径使得极限不存在或不一致即可.